\begin{abstract}
Continuous payments let a recipient spend an output before its creating
transaction executes publicly. Next is a native UTXO ledger in which a
recipient verifies a quorum-signed output certificate (TXCer) after one
certification round and can request a successor while public processing
continues. If the successor executes first, the committee records a funded
obligation against the direct issuer. Original signers can resubmit the
source from its publicly exposed certificate. A deadline-triggered decision
debits the issuer's reserve, and a late source repays it. Optional authorized
history commands then record the reserve reference at the original
coordinates while preserving owner authorizations, successor references,
and the complete BlockID. For four-member organizations and a four-validator
committee, each requiring three votes with at most one static Byzantine
member, we prove unique consumption, coverage of outstanding certificate
liability, and delay-neutral terminal CAL holdings for source-closed,
conflict-free payment sets. On one host with synchronous member commits,
a 100-hop chain reaches its last certified receipt in a median 4.557\,s,
versus 64.476\,s with per-hop waiting. Targeted boundary and restart tests
validate the conditions on capacity, economic closure, and historical reads.
\end{abstract}
\begin{IEEEkeywords}
Blockchain, certified payments, collateral, Byzantine fault tolerance,
chameleon hash.
\end{IEEEkeywords}
